% =====================================================================
%  REVISION of sec:gr-nonlocality in spacetime_from_non-computability.tex
%  Drafted 2026-06-13. Three reinforcing changes, all within this one
%  subsection:
%    (1) ER=EPR / Van Raamsdonk recast from "same direction of travel"
%        (convergence) to genuine CONTRAST (near-opposite mechanism):
%        they build a geometric bridge; we observe there was never a gap.
%    (2) The Vaxjo 2025 poster stated explicitly as the PRIORITY claim --
%        first public recognition that relational GR already accommodates
%        non-locality. Anchored to \autocite{Small2025} (existing key).
%    (3) The knowledge-space / co-location MECHANISM expanded (what JS
%        asked for): entangled systems share one knowledge-space location
%        and acquire distinct locations only by interacting with already-
%        located things. Includes the deeper-than-emergent-space subtlety
%        (co-location =/= coincidence) and the forward link to
%        sec:confinement.
%
%  HOW TO APPLY: this REPLACES the passage in sec:gr-nonlocality from
%  "The upshot is that general relativity, taken relationally and
%  seriously, already contains quantum non-locality..." (the sentence at
%  ~line 186) through the end of that paragraph at ~line 197
%  ("...flagged as a thesis rather than a theorem."). Everything BEFORE
%  that sentence (the sealed-bottle electron, the pre-spatial setup) is
%  retained unchanged; this revision picks up from the "upshot" and
%  extends it. The \medskip summary paragraph after (line 199+) is also
%  retained unchanged.
%
%  No new citations required (Small2025, VanRaamsdonk2010,
%  MaldacenaSusskind2013 all already present; fanout-boundary and Zurek
%  already cited elsewhere).
% =====================================================================

The upshot is a thesis about general relativity that the wider literature has not
drawn, and which the present author first stated publicly in the 2025 V\"axj\"o
poster \autocite{Small2025}: relational general relativity, taken seriously,
\emph{already} contains quantum non-locality. Entanglement is the pre-spatial
relational structure, and what looks like spooky action between distant points is
the resolution of a relation that was never spatially separated to begin with. The
claim is one of recognition rather than construction. Nothing need be added to
general relativity to accommodate the non-locality of quantum mechanics; one need
only notice that a theory in which spatial separation is fixed by the network of
interactions, and not by a pre-existing stage, has never committed to the
separation of two systems that have not yet interacted through that network. The
non-locality is already there, unrecognised, in the relational reading of the
theory, and the poster's contribution---advanced earlier in the author's CASYS'05
paper \autocite{small_2006} and developed here---was to point this out.

It is worth being exact about how this differs from the programmes it most
resembles, because the resemblance is real at the level of the phenomenon and
misleading at the level of the mechanism. Building spacetime from entanglement
\autocite{VanRaamsdonk2010} and the ER\,$=$\,EPR conjecture
\autocite{MaldacenaSusskind2013} share with the present account the conviction
that the spatial separation of entangled partners is not fundamental. They differ
from it---and from each other less than they differ from it---in what they take to
follow. Those programmes \emph{construct} a geometric connection: a wormhole, an
Einstein--Rosen bridge, a throat of new spatial geometry that realises the
entanglement as a connection through space. Entanglement, on that reading, builds
or is geometric structure linking the partners. The present account does the
opposite. It builds nothing and connects nothing, because there was never a gap to
span: the partners were never at two places, so no bridge between two places is
required. Where ER\,$=$\,EPR adds a geometric object to explain the correlation,
the relational reading subtracts an assumption---that the partners were
separated---and the puzzle dissolves rather than being mechanised. One approach
erects a structure across a gap; the other observes that relational general
relativity never opened the gap. This is not the same argument reached by a
different route, and it is not a corollary of the holographic results; it is a
distinct, and in its mechanism nearly opposite, account of the same fact, and it
is offered as the framework's interpretive thesis, flagged as a thesis rather than
a theorem.

The mechanism is made precise by reading the space of computations from the
observer's side, as a \emph{knowledge space}: a position in it is a state of
information, not a point of a manifold, and what fixes a system's position is what
is known of it through its interactions. Entanglement is co-location in this
space. Two systems in an entangled relation do not occupy two knowledge-space
positions awaiting a connection; they occupy \emph{one}, because the relation is a
single computational object that has not been resolved into parts. A system
acquires a \emph{distinct} location only by interacting with other systems that
already hold locations in the network---location is inherited relationally, a
place obtained by being related to things that have places. An entangled pair
shares one knowledge-space location until one partner is fanned out
(\S\ref{sec:fanout-boundary}) into the already-located network; that interaction
is what assigns the partner a distinct location, and the other partner's location
follows from the now-resolved relation. Spatial separation is in this sense
\emph{produced} by interaction, not revealed by it: the separation did not
pre-exist the measurement that established it.

One distinction must be drawn sharply, because without it the reading invites a
false consequence. Co-location in knowledge space is not coincidence in emergent
space. Knowledge space is the deeper, pre-spatial layer from which emergent
spacetime is built (\S\ref{sec:comp-distance}, \S\ref{sec:three-dimensions}); two
systems may share a knowledge-space location while the composite they belong to,
having interacted with the rest of the network, holds an emergent-space location
of finite extent. The point has a direct realisation in the strong sector,
developed in \S\ref{sec:confinement}. The three quarks of a colour singlet are
maximally entangled, so they share a single knowledge-space location; the
colour-singlet constraint, which forbids them the external entanglement that
location-acquisition requires, is exactly what prevents any one of them from
acquiring a distinct location of its own. Confinement is, on this reading,
pre-spatial co-location that cannot be resolved: a quark cannot be drawn to a
place apart because, in the knowledge space that is fundamental, it has no place
apart to be drawn to, and the singlet will not let it acquire one. That the proton
nonetheless has a finite radius in emergent space is no contradiction---the
emergent extent is the structure of the bound state as a whole, not a separation
of the singlet---and relating the knowledge-space co-location to the
emergent-space potential between the charges is precisely the open problem of
\S\ref{sec:confinement}.

The graded process by which new interactions create new knowledge-space locations,
and so create and sharpen emergent spacetime structure, is itself graded in the
manner of \S\ref{sec:fanout-boundary}---a location is as definite as the
interactions that fix it are redundant---and its dynamics is the subject of the
second paper. The present paper states what knowledge-space co-location is; how
locations proliferate, and how emergent distance is built from their proliferation,
is deferred.
